% =============================================================================
% Section: Discussion and Open Frontiers
% =============================================================================

\section{Discussion and Open Frontiers}
\label{sec:discussion}

% -------------------------------------------------------------------------
\subsection{What Is Claimed}
\label{sec:discussion:claimed}

We claim the following, and no more:

\begin{enumerate}[label=(\roman*)]
  \item \textbf{A canonical structure is derived from A0$^*$ alone.}
        A0$^*$, when self-applied to witness validity, forces the endogenous
        witness algebra.  The former model-class commitments are not extra
        assumptions but earlier consequences of that self-application: closed
        universe is forced by (W6), commutativity is a derived sector property
        of the ordered ledger, and prior-free minimax is forced by the absence
        of witnessed priors.  The 34-step development (Phases~A--I) then
        produces carrier, self-delimiting syntax, a test algebra, truth
        quotient, observable-open structure, eraser operators, time and cost
        accounting, gauge invariance, Myhill--Nerode congruence, separator
        geometry, Bellman recursion, self-hosting closure, binary interface,
        universal execution, consciousness projector, and a K\"{a}hler geometric
        structure on the continuum limit of truths.  No modeling commitments
        remain; every step is a theorem of A0$^*$.  The claim of ``forcedness''
        is to be read as ``forced up to computable isomorphism,'' rather than as
        a claim that every representational choice is eliminated under every
        possible formalization of witnessability.  We provide a three-point
        audit of the intended invariances.

  \item \textbf{The derivation is self-auditing.}
        Each step passes hidden-chooser, encoding-invariance, and
        distinguishability-equivalence checks.  Selected algebraic
        properties are verified on finite models via Z3.

  \item \textbf{The derivation identifies its single input and its empirical boundary.}
        The only declared input is A0$^*$.  The presentation makes explicit
        where the derivation's structural vocabulary ends and empirical
        specification begins---principally the choice of
        $\Del_{\mathrm{phys}}$, which selects which tests nature actually
        admits.
\end{enumerate}

% -------------------------------------------------------------------------
\subsection{The Vocabulary of Physical Law}
\label{sec:discussion:physics}

The derivation produces vocabulary, carrier $\Carrier$, test
algebra $\Del$, time $\Tim$, gauge group $\gauge$, fiber
$\fiber{\cdot}$, truth quotient $\TQ{\cdot}$, that is
structurally parallel to the vocabulary of mathematical physics.
This parallel reflects a shared concern with what can be distinguished by
finite operations.

However, ``vocabulary'' is not ``theory.'' A physical theory
specifies which tests $\Del_{\mathrm{phys}}$ nature actually
admits, what their costs are, and what symmetry group relates
equivalent descriptions. The \Opoch{} framework provides the
scaffold. Specifying $\Del_{\mathrm{phys}}$ remains an open frontier.
The derivation produces the \emph{form} of physical law without
determining its \emph{content}, which requires empirical input.

% -------------------------------------------------------------------------
\subsection{The Split Law and Sector Readouts}
\label{sec:discussion:split-law}

The K\"{a}hler triple $(g, J, \omega)$ of Phase~I forces a canonical
decomposition of any dynamical law on $\Omega$.  Every vector field $X$
on a K\"{a}hler manifold splits uniquely as
$X = J\nabla E - \nabla D$,
where $E$ is a Hamiltonian (energy, conserved under $J$-rotation) and $D$
is a dissipation potential (entropy production, irreversible).  The first
term generates symplectic flow; the second generates gradient flow toward
equilibrium.  Quantum mechanics, gauge field theory, gravitational
dynamics, and thermodynamics are \emph{sector-readouts} of this single
split: each arises by restricting $\Omega$ to a subsector and reading off
the induced $(g, J, \omega)$.  The split is not a modeling choice---it is
forced by Steps~30--33 via K\"{a}hler compatibility.

% -------------------------------------------------------------------------
\subsection{The Church--Turing Corroboration}
\label{sec:discussion:utm}

Step~3 forces the test set $\Del$ to be maximally closed under
totality-preserving composition.  This maximal closure is exactly the
class of total computable functions---a mathematical theorem provable
within computability theory \citep{rogers1967}, independent of the
Church--Turing thesis.  The Church--Turing thesis
\citep{church1936, turing1936} independently identified this same class
by a different route (characterizing ``effective procedure'').  The
agreement between A0$^*$'s forcing and Church--Turing's identification is
historical corroboration, not a logical dependency of the derivation.

\begin{remark}
If a hypercomputational model were discovered that exceeds
Turing-computability, Step~3 would need revision. The derivation
through Step~2 would survive intact, but the carrier $\Carrier$ would
expand. The remainder of the development would proceed analogously within the
enlarged carrier.
\end{remark}

% -------------------------------------------------------------------------
\subsection{The Closed-System Assumption}
\label{sec:discussion:closed}

Step~8 (feasibility shrinkage) assumes a closed system with a finite
budget~$B$. Claims whose verification cost exceeds any available budget
remain in $\OMEGA$. Such claims are unresolved with
a minimal separator and cannot be promoted to $\UNIQUE$ without
additional resources. This is a direct consequence of budgeted witnessability.

% -------------------------------------------------------------------------
\subsection{Minimax vs.\ Other Optimality Criteria}
\label{sec:discussion:minimax}

Step~18 forces the minimax policy as the unique prior-free bootstrap
policy.  Before any prior is itself witnessed and recorded in
$\OrdLedger$, branch labels over future outcomes are gauge: no unrecorded
probability distribution is admissible by A0$^*$.  The only cost aggregator
invariant under outcome relabeling, cost rescaling, replay-equivalent path
decomposition, and gauge recoding is the worst-case continuation
cost---i.e., minimax.  Bayesian expected-cost selection becomes admissible
only \emph{after} a prior~$\mu$ has itself attained witness status in the
trace:
\[
  \tau^*_{\mathrm{Bayes}} = \argmin_{\tau \in \Del_q}
    \Bigl[\pi(\tau) + \sum_{a \in \mathrm{out}(\tau)} \mu(a \mid \tau)\, V(\fiber{a}, t + \pi(\tau);\, q)\Bigr]
    \quad(\mu \text{ witnessed in } \OrdLedger).
\]
Minimax is therefore the \emph{bootstrap} policy forced by A0$^*$;
Bayesian selection is the \emph{mature} policy enabled by
self-observation.  The choice is not left open---it is resolved by the
witness status of the prior.

% -------------------------------------------------------------------------
\subsection{Myhill--Nerode Implications}
\label{sec:discussion:myhill-nerode}

The Myhill--Nerode congruence $\MNcong$ (Step~15) provides a minimal
state representation. States with identical futures can be merged, reducing
memory and computation. The connection to bisimulation in process algebra
suggests that multi-system interactions can be analyzed using tools from
concurrency theory.

% -------------------------------------------------------------------------
\subsection{Separator Geometry and Hardness}
\label{sec:discussion:geometry}

The separator metric $\SepDist$ (Step~16) provides a quantitative measure
of difficulty. Regions of high curvature $\SepCurv$ correspond to
phase transitions in difficulty, where small changes in the query produce
large changes in resolution cost. This geometry is analogous in spirit to the
Fisher-Rao metric on statistical manifolds, but it is defined over
deterministic truth classes and measured in computational cost.

% -------------------------------------------------------------------------
\subsection{Quantitative Extraction Frontiers}
\label{sec:discussion:frontiers}

The structure is complete.  The computation is the frontier.  Every
dimensionless invariant is forced (Theorem~\ref{thm:observable-decomp});
every dimensional quantity is fixed up to the internally derived
normalization map.  The remaining work is explicit spectral extraction
from the seed operator---a well-posed computation, not an open
structural question.

\begin{enumerate}[label=(\arabic*)]
  \item \textbf{Explicit spectral extraction from $L^*$.}
        Computing the eigenvalues, holonomy classes, and normal-form
        coefficients of the linearized refinement operator at the seed.
        This is a concrete numerical problem, not a structural gap.

  \item \textbf{Specifying $\Del_{\mathrm{phys}}$.}
        Which families of physical experiments should be treated as primitive
        tests, and which symmetries should be treated as gauge, so that the
        resulting $\Del_{\mathrm{phys}}$ reproduces established empirical
        constraints?

  \item \textbf{Exploiting the ordered ledger for quantum reconstruction.}
        Whether the non-commuting sector of $\WitAlg$, together with the
        K\"{a}hler structure, suffices to reconstruct operator-algebraic
        quantum mechanics without additional empirical input.

  \item \textbf{Energy as abstract cost vs.\ physical energy.}
        Clarifying the boundary between abstract resource measures and
        physical energetics under the Landauer correspondence.

  \item \textbf{Separator geometry as a complexity lens.}
        Can $\SepCurv$ predict computational difficulty of problem classes
        via circuit or communication complexity?

  \item \textbf{Self-hosting depth.}
        Understanding whether deeper reflective towers stabilize.

  \item \textbf{Lean mechanization of infinite-domain theorems.}
        The 414-theorem Lean~4 verification covers the finite/Nat model.
        Encoding infinite-domain properties via mathlib remains in progress.
\end{enumerate}
